\chapter{Phase-boundary sets}
\label{ch:phasesets}

\section{Why a separate concept}

The volume of a closed compact body is the familiar surface integral. For the other surfaces
of \mftl{} the question ``what volume does this enclose?'' may have no answer, or an answer that
depends on other surfaces:
\begin{itemize}
  \item a single flat film that wraps $x$ and $y$ of a periodic box separates nothing --- the
    region above it is connected to the region below it through the periodic $z$ faces;
  \item two such films together bound a slab, whose volume is well defined;
  \item an open cap sitting on a wall orifice bounds a phase only together with the wall;
  \item a bubble in a fully periodic box bounds two phases: the bubble, and the matrix around it.
\end{itemize}
A \emph{phase-boundary set} makes the question explicit: it lists the surfaces that together
bound one phase, with an orientation for each, plus optional wall closures, and every
closed-phase query on it returns a status alongside its value.

\section{Status model}

\begin{lstlisting}[numbers=none]
typedef enum { MFTL_DEFINED = 0, MFTL_UNDEFINED = 1, MFTL_UNSUPPORTED = 2 } MFTL_PhaseStatus;
\end{lstlisting}
\begin{description}
  \item[\texttt{MFTL\_DEFINED}] the quantity is meaningful and has been evaluated;
  \item[\texttt{MFTL\_UNDEFINED}] the topology does not define it (the set does not separate
    space, its orientations are inconsistent, an open member is not closed by a declared wall, or
    the region is empty or fills the box);
  \item[\texttt{MFTL\_UNSUPPORTED}] it is mathematically defined, but this implementation cannot
    evaluate it (for example, no admissible evaluation axis).
\end{description}
No query returns a value with a status other than \texttt{MFTL\_DEFINED}; the output argument is
left untouched then. Quantities that depend only on the surface --- area, surface centre, boundary
loops --- are always defined.

\section{Building a set}

\begin{lstlisting}[numbers=none]
MFTL_PhaseSet *MFTLPhaseSetCreate(MFTL_Domain *domain);
int  MFTLPhaseSetAddMember(MFTL_PhaseSet *S, MFTL_Object *body, int sigma);
int  MFTLPhaseSetAddWallClosure(MFTL_PhaseSet *S, int axis, int side); /* MFTL_WALL_MIN/MAX */
void MFTLPhaseSetDestroy(MFTL_PhaseSet **S);
\end{lstlisting}
A member contributes all its connected components with the orientation sign
$\sigma=\pm1$: with $\sigma=+1$ the phase lies on the side into which the stored normals do
\emph{not} point (inside an outward-wound closed body). A wall closure declares that the phase
may be closed by the wall at the lower or upper end of a \emph{non-periodic} axis. Members are
assumed to be embedded (no self- or mutual intersections); orientability and manifoldness are
checked by the builders.

\section{When a set bounds a phase}

Let $\hat S$ be the members together with the wall patches that close their boundary loops. A
set bounds a finite phase region if and only if
\begin{enumerate}
  \item $\hat S$ is closed: every boundary loop of every member lies on a declared closure wall
    (within a tolerance), and the wall patches can be chosen so that the loops are closed;
  \item its net flux vanishes on every axis:
    \begin{equation}
      \Phi_j(\hat S) = \sum_m \sigma_m \oint_m n_j\,dA + \sum_W \Phi_j(P_W) = 0,\qquad j=x,y,z;
      \label{eq:flux}
    \end{equation}
  \item the orientation is consistent: walking through space, one crosses the members
    alternately into and out of the phase.
\end{enumerate}
A single wrapping film violates (\ref{eq:flux}): its flux along $z$ is $\pm$ the area of the box
face. Two films with opposite normals (lower film pointing down, upper film pointing up, both
$\sigma=+1$) have zero total flux and bound the slab between them.

\section{Evaluation}

\paragraph{Evaluation axis.} The formulas integrate along one \emph{admissible} axis $k$: the
wall-closure axis if there is one, otherwise a non-periodic axis, otherwise a periodic axis along
which every member is compact (ties go to the smallest index). Without an admissible axis the
status is \texttt{MFTL\_UNSUPPORTED} (for example two tilted planes that wrap all three axes).

\paragraph{Volume.} With continuous coordinates $\vect{Y}$ relative to the set's reference point,
\begin{equation}
  V_{\text{raw}} = \sum_m \sigma_m \oint_m Y_k\, n_k\,dA + \text{(wall-patch terms)}.
  \label{eq:vraw}
\end{equation}
On a closed compact body this is the divergence theorem, $V=\oint Y_k n_k\,dA$. On a periodic
axis the connected components are classified as \emph{compact} (zero flux along $k$) or
\emph{layers} (flux $\pm A_k$, $A_k$ being the box face area normal to $k$), and
\begin{equation}
  V = \begin{cases}
    V_{\text{raw}} \bmod V_{\text{cell}} & \text{periodic $k$ with layers,}\\
    V_{\text{raw}} + \lambda_E V_{\text{cell}} & \text{periodic $k$ without layers,}\\
    V_{\text{raw}} & \text{walled $k$,}
  \end{cases}
\end{equation}
where $V_{\text{cell}}=A_k L_k$ and $\lambda_E\in\{0,1\}$ is the phase level of the exterior (1 for
the matrix phase around a bubble with $\sigma=-1$). A volume within a small tolerance of $0$ or of
the box volume is \texttt{MFTL\_UNDEFINED} (empty or full region).

\paragraph{Centre and inside/outside.} The centre along the evaluation axis follows from
$\sum_m\sigma_m\oint (Y_k+\pi_c)^2 n_k\,dA/(2V)$ with a canonical placement $\pi_c$ of the
components, and is reported modulo the period on a periodic axis. Along a periodic axis that the
phase region wraps, the centre is \texttt{MFTL\_UNDEFINED}; for a region with several components
along a periodic axis it is \texttt{MFTL\_UNSUPPORTED}. \fn{MFTLPhaseSetInside} casts a line
along the evaluation axis through the query point and inspects the nearest crossing below it.

\section{Queries}

\begin{lstlisting}[numbers=none]
MFTL_PhaseStatus MFTLPhaseSetClassify(MFTL_PhaseSet *S, MFTL_PhaseSetInfo *info);
MFTL_PhaseStatus MFTLPhaseSetVolume(MFTL_PhaseSet *S, double *V);
MFTL_PhaseStatus MFTLPhaseSetCentre(MFTL_PhaseSet *S, double C[3], MFTL_PhaseStatus status[3]);
MFTL_PhaseStatus MFTLPhaseSetInside(MFTL_PhaseSet *S, const double Q[3], int *where);
MFTL_PhaseStatus MFTLPhaseSetVolumeRecovery(MFTL_PhaseSet *S, MFTL_Ensemble *ens,
                                            double V_target, double *V_out);
\end{lstlisting}
All are collective: the triangles of the set are gathered on every rank and evaluated
identically. \fn{MFTLPhaseSetClassify} fills an \texttt{MFTL\_PhaseSetInfo} with the status, a
human-readable reason, the flux vector, the evaluation axis and the component counts.
\fn{MFTLPhaseSetInside} sets \texttt{*where} to \texttt{MFTL\_PHASE\_INSIDE},
\texttt{MFTL\_PHASE\_OUTSIDE} or \texttt{MFTL\_PHASE\_ON\_BOUNDARY}. \fn{MFTLPhaseSetVolumeRecovery}
moves every member by a uniform displacement along its phase-outward vertex normal
($\sigma\times$ the stored normal; boundary vertices on wall closures stay fixed) by
$\delta=(V_{\text{target}}-V)/A(S)$, iterated (at most 8 steps) until
$|V-V_{\text{target}}|\le10^{-12}V_{\text{cell}}$. The tolerances of the evaluation can be read and
changed with \fn{MFTLPhaseSetGetTolerances} / \fn{MFTLPhaseSetSetTolerances}.

\section{Summary of outcomes}

\begin{center}\small
\begin{tabularx}{\textwidth}{@{}Xl@{}}
\toprule
Situation & Status \\
\midrule
closed compact body, $\sigma=+1$ & DEFINED (its own volume) \\
bubble with $\sigma=-1$ in a fully periodic box & DEFINED ($V_{\text{cell}}-V_{\text{bubble}}$) \\
single film wrapping $x$ and $y$ & UNDEFINED ($\Phi_z\neq0$) \\
two wrapping films with opposite normals & DEFINED (the slab) \\
film with normal $+z$ over a $z$-min wall closure & DEFINED (the layer under the film) \\
cap on a wall orifice with a wall closure & DEFINED (cap plus disk) \\
open surface without a declared closure & UNDEFINED \\
no admissible axis (tilted triply-wrapping planes, diagonal winding) & UNSUPPORTED \\
loop that winds around a periodic direction of the wall; loops on two walls; ambiguous closure & UNSUPPORTED \\
\bottomrule
\end{tabularx}
\end{center}

The example of \S\ref{sec:ex-phaseset} evaluates the single-film and two-film cases.
